\chapter{مقدمة في شرح الرسالة وخطبة الإمام الشافعي}

\section{مقدمة الدرس وخطبة الحاجة}
إن الحمد لله نحمده ونستعينه ونستغفره، ونعوذ بالله تعالى من شرور أنفسنا ومن سيئات أعمالنا، من يهده الله فلا مضل له، ومن يضلل فلا هادي له. وأشهد أن لا إله إلا الله وحده لا شريك له، وأشهد أن محمداً عبده ورسوله.

﴿يَا أَيُّهَا الَّذِينَ آمَنُوا اتَّقُوا اللَّهَ حَقَّ تُقَاتِهِ وَلَا تَمُوتُنَّ إِلَّا وَأَنتُم مُّسْلِمُونَ﴾، ﴿يَا أَيُّهَا النَّاسُ اتَّقُوا رَبَّكُمُ الَّذِي خَلَقَكُم مِّن نَّفْسٍ وَاحِدَةٍ وَخَلَقَ مِنْهَا زَوْجَهَا وَبَثَّ مِنْهُمَا رِجَالًا كَثِيرًا وَنِسَاءً وَاتَّقُوا اللَّهَ الَّذِي تَسَاءَلُونَ بِهِ وَالْأَرْحَامَ إِنَّ اللَّهَ كَانَ عَلَيْكُمْ رَقِيبًا﴾، ﴿يَا أَيُّهَا الَّذِينَ آمَنُوا اتَّقُوا اللَّهَ وَقُولُوا قَوْلًا سَدِيدًا * يُصْلِحْ لَكُمْ أَعْمَالَكُمْ وَيَغْفِرْ لَكُمْ ذُنُوبَكُمْ وَمَن يُطِعِ اللَّهَ وَرَسُولَهُ فَقَدْ فَازَ فَوْزًا عَظِيمًا﴾.

أما بعد؛ فإن أصدق الحديث كتاب الله تعالى، وإن خير الهدي هدي محمد عليه الصلاة والسلام، وإن شر الأمور محدثاتها، وإن كل محدثة بدعة، وكل بدعة ضلالة، وكل ضلالة في النار.

هذا هو الدرس الأول في القراءة في كتاب «الرسالة» للإمام أبي عبد الله الشافعي رحمه الله تعالى. وقد اهتم الشافعي رحمه الله في كتابه بأمرين عظيمين: الأول: تحرير القواعد الأصولية، وإقامة الأدلة عليها من الكتاب والسنة، وتأييد كلامه بشواهد من اللغة العربية. الثاني: الإكثار من الأمثلة لزيادة الإيضاح والتطبيق لقضايا أصول الشريعة وفروعها، مع مناقشة المخالفين بأسلوب سهل جزل. وهو أول ما دُوّن في أصول الفقه، ويكفينا شرفاً أن نتلقى عن الإمام الشافعي، نقرأ له ونتعلم منه، ونتأدب بأدبه، ونحفظ شيئاً من القواعد التي أصّلها.

\section{أهمية كتاب الرسالة وتحقيقاته}
ألّف الشافعي «الرسالة» مرتين؛ الرسالة القديمة التي أرسلها لعبد الرحمن بن مهدي حين سأله عن دلالات النصوص والإجماع والقياس، ثم أملى الرسالة الجديدة على تلميذه الربيع بن سليمان المرادي.

يقول أبو القاسم عثمان بن سعيد الأنماطي عن الإمام المزني رحمه الله: «أنا أنظر في كتاب الرسالة عن الشافعي منذ خمسين سنة، ما أعلم أني نظرت فيه مرة إلا وأنا أستفيد منه شيئاً لم أكن عرفته». فظل خمسين سنة يراجع ويذاكر هذا الكتاب الجليل.

وقد حُقق هذا الكتاب تحقيقاً جليلاً على يد شيخ شيوخنا أبي الأشبال أحمد بن محمد شاكر رحمه الله، على نسخة الربيع بن سليمان، وهي من أقدم المخطوطات على وجه الأرض (كُتبت سنة 265هـ). كما حققها الدكتور رفعت فوزي عبد المطلب عند تحقيقه لموسوعة «الأم».

\section{قراءة إسناد الكتاب ونسب الإمام الشافعي}
قال الربيع بن سليمان: بسم الله الرحمن الرحيم. أخبرنا أبو عبد الله محمد بن إدريس بن العباس بن عثمان بن شافع بن السائب بن عبيد بن عبد يزيد بن هاشم بن المطلب بن عبد مناف.

والمطلب بن عبد مناف كان أخاً لهاشم بن عبد مناف، والمطلب هو الذي ذهب إلى المدينة ليأتي بابن أخيه هاشم (شيبة الحمد)، فأردفه خلفه ودخل مكة، فقال القرشيون: «من هذا؟»، قالوا: «هذا المطلب»، وقالوا عمن معه: «عبد المطلب» يظنونه عبداً اشتراه، فصار «عبد المطلب» لقباً اشتهر به جد النبي صلى الله عليه وسلم. والشافعي قرشي صليبة من بني المطلب، أبناء عمومة رسول الله صلى الله عليه وسلم.

\section{شرح خطبة الإمام الشافعي في الرسالة}
ابتدأ الشافعي كتابه بالبسملة والحمد، اقتداءً بالكتاب الكريم. قال رحمه الله:
«الحمد لله الذي خلق السماوات والأرض وجعل الظلمات والنور ثم الذين كفروا بربهم يعدلون. والحمد لله الذي لا يُؤدّى شُكر نعمة من نعمه إلا بنعمة منه تُوجب على مُؤدّي ماضي نعمه بأدائها نعمة حادثة يجب عليه شكره بها...»

\subsection{الفرق بين الحمد والشكر}
الحمد يكون باللسان والقلب على جهة التعظيم، ويكون على النعمة وغيرها. أما الشكر فلا يكون إلا على نعمة أو جميل يُقدّم، ويكون بالقلب واللسان والجوارح. فالحمد أعم سبباً (لأنه على النعمة وغيرها)، والشكر أعم متعلقاً (لأنه بالجوارح والقلب واللسان). ونحن نحمد الله على كمال صفاته وعلى عظيم إنعامه.

\subsection{فقه القلوب عند الشافعي}
في قول الشافعي: «لا يُؤدّى شكر نعمة من نعمه إلا بنعمة منه...» درس تربوي عظيم؛ فالعبد لا يستطيع أداء شكر الله إلا بإعانة الله وتوفيقه، وهذا التوفيق نعمة جديدة توجب شكراً جديداً. فالمؤمن دائم التقلب في نعم الله، ودائم الشكر لها، وكلما شكر زادته النعم، ﴿لَئِن شَكَرْتُمْ لَأَزِيدَنَّكُمْ﴾.

ثم قال: «ولا يبلغ الواصفون كُنه عظمته، الذي هو كما وصف نفسه وفوق ما يصفه به خلقه». وهذا أصل عقدي، فلا يجوز وصف الله بغير ما وصف به نفسه، ولا يجوز التكييف أو التمثيل، لأن الخلق لن يبلغوا حقيقة عظمته سبحانه.

\section{حال الخلق عند بعثة النبي صلى الله عليه وسلم}
ذكر الشافعي انقسام الناس عند مبعث النبي صلى الله عليه وسلم إلى قسمين:
الأول: أهل الكتاب (اليهود والنصارى) الذين بدلوا أحكام كتبهم وحرّفوها.
الثاني: المشركون عبدة الأصنام والأوثان من العرب، وعبدة النيران والنجوم من الأعاجم. فبعث الله محمداً صلى الله عليه وسلم لإنقاذ البشرية.

وفي قوله: «فلما بلغ الكتاب أجله فحقّ قضاء الله بإظهار دينه الذي اصطفى بعد استعلاء معصيته التي لم يرضَ...» لمحة تربوية؛ فالتمكين وإظهار الدين يكون بقضاء الله، والواجب علينا الأخذ بالأسباب الشرعية وعدم الركون لمعصية الله طمعاً في التمكين.

\section{خصائص النبي محمد صلى الله عليه وسلم ومقاصد رسالته}
اصطفى الله محمداً صلى الله عليه وسلم من خيار خيار العرب، وختم به النبوة، وعمّ برسالته الناس كافة. وخصّ قومه (العرب) بالنذارة أولاً ﴿وَأَنذِرْ عَشِيرَتَكَ الْأَقْرَبِينَ﴾ ليكونوا حملة الرسالة إلى غيرهم.

\subsection{الاعتزاز باللغة العربية ومقاصده}
يؤكد الشافعي أن القرآن نزل بلسان عربي مبين خالص لا تشوبه عُجمة. والاعتزاز بالعربية هو اعتزاز بالدين؛ فمتابعة الهدي الظاهر في اللغة والسمت تورث الاهتداء في الباطن. وقد نقل الشيخ أحمد شاكر عن والده محمود شاكر تحذيره من دعوات الاستعمار لترجمة القرآن وجعله أعجمياً في مستعمراتهم، لإذابة الشخصية الإسلامية وقطع الرابطة العربية بين المسلمين.

\section{آداب طالب العلم ومسلك النجاة}
ختم الشافعي هذه المقدمة ببيان طبقات الناس في العلم، وأن العلم الحقيقي هو علم الكتاب والسنة. وأوصى طلبة العلم بأمور:
1. بلوغ غاية الجهد في الاستكثار من العلم.
2. الصبر على العوارض (الفقر، التعب، الانشغال).
3. إخلاص النية لله عز وجل.
4. الرغبة إلى الله والتذلل له لطلب العون والفهم.

فمن أدرك أحكام الله نصاً واستدلالاً، ووفقه الله للقول والعمل، فاز بالفضيلة في دينه ودنياه، ونوّرت في قلبه الحكمة.